\section{Conclusion}
\label{sec:conclusion}

A frozen Stable Diffusion 1.5 model adapted from 50 coordinate-corresponding pairs can be made to normalise
between two scanners, but only after replacing the planned recipe. A target-only colour LoRA with
ControlNet and LCM barely moves colour. Source conditioning, DDIM inversion and a post-hoc fusion that
returns the learned colour change to the source pixels reach SSIM 0.729 and beat the classical baselines on
typical-slide colour. The VAE bottleneck explains most of the gap to the raw scan, and fine-tuning its
decoder raises SSIM mainly by smoothing. Learned nondiffusion baselines trained on the same 50 pairs
remain better on colour and SSIM. On nucleus preservation, ParamNet and the fused pipeline both pass the
Relative Dice test (0.999 and 0.994), and the unfused diffusion output does not. For
few-shot scanner normalisation, a well-chosen colour map currently beats a resynthesising generative
model; diffusion is most useful here as a source of colour change that fusion projects back onto the
original image.

\section*{Acknowledgment}
Computations were performed on the University of the Witwatersrand's mscluster
high-performance computing cluster. We thank the organisers of
MITOS-ATYPIA-14, CAMELYON17, Lizard, PanNuke and TCGA for making their data public.
